\section{`Two goods' models \label{sec:twogoodsmodels}}

\begin{comment}
%%%%%%%%%%%%%%%%%%%%%%%%%%%%%
% bits for this section
%%%%%%%%%%%%%%%%%%%%%%%%%%%%%
WUD 3	Model: is it possible that stock increasing forever is a good strategy for a monopolist in this model? i.e. it's better for firms to produce more even if they burn the produce? Or are firms just getting it wrong? Note: increasing revenue with one firm just = increasing money velocity. Does it really mean anything?


\end{comment}

\subsection{How should agents optimise? \label{sec:2goodsOptimise}}

%The point of this section: in actually working out how agents should optimise, a whole series of interesting things occurred. I want to write up those issues here. So, what are they?

%I will want to have already laid out the structure of twogoods1 above, so I can assume that info has already been covered. Also, somewhere, I'm going to need a more detailed discussion of the nature of time in these models - it's a key problem. I have to mention that here...

\cmt{crossref to discussions on merits of utility and problems of optimising in agents.}


%Secondly: at each stage of modelling, the choices that present themselves can take the model in many directions. It is important, therefore, to remember what the ultimate goal is, and to try and judge when a particular choice is relevant to that ultimate goal. So the principle here is: does the discussion help illuminate the differences between analytic and agent modelling, or does it further the goal of understanding the impact of space cost changes? If it does neither of those two, it should be avoided if possible.
%\cmt{That's surely an instruction for me, not a PhD para..?}

%Things to cover here: 
%Price, difficulty of numeraire approach
%Consequences of full-information optimisation; advantages of blind search e.g one-day optimising people means a catch-22: no stock = no point in working; no work = no stock. Does that mean agents aren't 'rational?' The more blindly they can find their way to an optimum, the more it can be said to objectively exist?
%Brute force vs analytic optimisation




%This is, however, not nearly as easy as one might hope as the aforementioned conditions of uncertainty preclude a simple formulation. As described below, 
%Moved this to Fowler's own section:
%%---
%\begin{quote}
%``The logical step to improve the functionality of this model is to develop equations that account for the rationing of wages and production, based on worker employment (for wages) and availability (for production)... the model must be reformulated from its most basic assumptions before such rationing can be incorporated to move forward in this highly desirable direction.'' (\citealt{fowler_taking_2007} p.277)
%\end{quote}

%This is, of course, exactly the kind of argument Hayek was attempting to counter: success in providing a model rationing system would end up saying nothing about the fundamental problem of how disaggregated actors can coordinate production. There are benefits, however, to the `subjective' approach of testing possible goods bundles against utility, parallel to the benefits of blind search: no prior knowledge of the state of the market is required. Actors can act based on the informatin presented to them.
%%---

%Note first that the `budget' for the brute force approach and the first constrained optima is, in fact, just the total quantities of good; price does not enter into the calculation. In the core model, price is analysed away from the numeraire good - but it is still meant to be part of the analysis. For the agent model, price cannot so easily be included, even implicitly, in utility calculations. 

%If money is involved, before it can give a value, that money needs converting into good amounts. That means it is impossible, given a certain wage and certain good prices, to know directly what the optimal spend would be.


%Ha: OK, the obvious reason why #3 won't work either: there's no guaranteed single price, and you can't sum Engel-function utilities.





\begin{comment}
CHUNKS TO WORK FROM

And here? Well, let's just take one aspect: decision timings. Firms 'decide' only once for every person having decided. So, all people could well have put their capital into one firm before another realises its price is incorrect. The raw number of people, then, will actually make a difference to the outcome of the firm, relative to other variables like how much capital is available. This makes finding a consistent way to model complicated: but I need to find a way to consistently describe it, and to explore the spaces created. This is where I'm having problems - perhaps simply because I feel overwhelmed. 

I've tried limiting the number of actors, and that's helping to get closer to what the dynamics are, but it's still proving vexatious. For example: if I increase the granularity of the utility parameter space search, firms adjust prices more slowly. Why? Somewhere in there, there's an intuitive answer: when people can make more fine-grained distinctions, it leads to more fine-grained adjustments in response. Seems obvious enough. But how to actually show that? And do I need to? I should, really.

Well, having just run another Torus model, I realise that space - introducing differentiation as it does - solves many of the problems inherent in point-economies, and does away with the analytic need for a product differentiation function like the CES. (Nice sentence, might steal that!) 

Indeed, Jacobs' "data / meaning of data / appropriate responses" thing - I'm actually using as an analysis tool. 

The problem in more detail: an explicit numeraire doesn't seem to work. Whatever good is acting as numeraire, in an agent model, actually has to get used. So say good1 is numeraire, and is produced. How do good2 producers use it to pay for anything?
Let's make a few numeraire notes:
�	So, Brakman - we know all this. The price of food is always equal to one. Other goods expressed relative to that. What does that look like mathematically? For example, how would one describe the price of another M variety, relative to the food price?

Note, if I do any groping here, it's going to be 'live groping' (brilliant!) - that is, each grope will have a real effect on the market. Goods and money will be exchanged - there's no 'groping' to find a set of clearing prices in this situation.

There's a lot to be written up here about the relationship between randomness, `optimality', groping, perfect knowledge. In some cases, perfect knowledge makes the model less likely to work...

Firm: 
I have the money - I'll give it to you in exchange for your time.
Person:
If I get that money, how much of other goods can I exchange it for?

\end{comment}